% Part: proof-theory
% Chapter: proof-search
% Section: completeness

\documentclass[../../../include/open-logic-section]{subfiles}

\begin{document}

\olfileid[hi]{pt}{ps}{cpl}

\olsection{पूर्णता}

अब हम दिखाते हैं कि प्रमाण-खोज कलन-विधि अचूक है; अर्थात् यदि वह बीच में रुक
जाती है या कभी समाप्त नहीं होती, तो आरंभिक अनुक्रम~$\Gamma \Sequent \Delta$
का कोई !!{proof} नहीं है। इसके लिए हम ऐसी !!a{structure}~$\Struct M$
परिभाषित करते हैं कि प्रत्येक $!A \in \Gamma$ के लिए $\Sat{M}{!A}$ और
प्रत्येक~$! \in \Delta$ के लिए $\Sat/{M}{!A}$ हो। दूसरे शब्दों में,
\Struct M में~$\Gamma$ के सभी !!{formula}s सत्य और~$\Delta$ के सभी
!!{formula}s असत्य हैं; अर्थात् $\Gamma \Sequent \Delta$ वैध नहीं है। चूँकि
\Log{G3c} निर्दोष है, $\Gamma \Sequent \Delta$ का कोई !!{proof} नहीं है।

हम !!{formula}s के युग्म $\Theta, \Xi$ और !!{constant}s के समुच्चय~$C$ से
$\Struct{M}$ को परिभाषित करते हैं। $\Theta, \Xi$ और~$C$ को प्रमाण-खोज के
विफल (बीच में रुके या कभी समाप्त न होने वाले) निष्पादन की एक \emph{विफलता
शाखा} से प्राप्त करते हैं। विफलता शाखा अनुक्रमों $\Pi_n \Sequent \Lambda_n$
और !!{constant}s के समुच्चयों~$C_n$ का ऐसा क्रम है कि $\Pi_0 \Sequent
\Lambda_0$ अंतिम अनुक्रम $\Gamma \Sequent \Delta$ है, और $\Pi_{n+1}
\Sequent \Lambda_{n+1}$, चरण~$n+1$ में $\Pi_n \Sequent \Lambda_n$ से उत्पन्न
ऐसा आधार-वाक्य है जिसका कलन-विधि कोई !!a{proof} नहीं खोजती। प्रत्येक~$n$ के
लिए ऐसा आधार-वाक्य स्पष्टतः होना चाहिए, अन्यथा कलन-विधि !!a{proof} खोज चुकी
होती। $C_n$ को~$\Pi_n \Sequent \Delta_n$ से संबद्ध !!{constant}s का समुच्चय
मानते हैं।

यदि कलन-विधि केवल परमाण्विक !!{formula}s वाले किसी सर्वोच्च अनुक्रम पर बीच में
रुकती है, तो उस अनुक्रम पर समाप्त होने वाली शाखा विफलता शाखा है। यदि
कलन-विधि कभी समाप्त नहीं होती, तो वह उत्तरोत्तर बड़े वृक्ष उत्पन्न करती रहती
है। इन्हें किसी अनंत वृक्ष को बढ़ाते जाने के रूप में समझा जा सकता है (वह वृक्ष
जो कलन-विधि के अनंत चरण चलने पर मिलता)। यह वृक्ष अनंत, किंतु सीमित-शाखी होगा
(प्रत्येक आसंधि की केवल एक या दो संतानें—उसके ठीक ऊपर के अनुक्रम—हैं)। क्योनिग
के प्रमेय से प्रत्येक अनंत सीमित-शाखी वृक्ष में कोई अनंत शाखा अवश्य होती है।
खोज कलन-विधि के अनंत काल तक चलने की स्थिति में ऐसी अनंत शाखा विफलता शाखा
होगी।

यदि $\Pi_n \Sequent \Lambda_n$ और $C_n$ का क्रम विफलता शाखा है, तो $\Theta
= \bigcup_n \Pi_n$ और $\Xi = \bigcup_n \Lambda_n$ (सूचकांक और !!{formula}s
की अनेक आवृत्तियाँ हटाकर), तथा $C = \bigcup_n C_n$ लें।

अब हम~$\Struct{M}$ को परिभाषित कर सकते हैं।
\begin{enumerate}
\item प्रांत $\Domain{M}$ उन पदों का समुच्चय है जिन्हें~$C$ और $\Gamma
\Sequent \Delta$ के !!{function}s से बनाया जा सकता है, यदि कोई हों, किंतु
बिना !!{variable}s के।
\item यदि $c \in \Domain{M}$, तो $\Assign{c}{M} = c$।
\item यदि $f$, $\Gamma \Sequent \Delta$ में कोई $m$-स्थानी !!{function} है और
$t_1$, \dots, $t_m \in \Domain{M}$, तो $\Assign{f}{M}(t_1, \dots, t_m) =
f(t_1, \dots, t_n)$।
\item यदि $R$, $\Gamma \Sequent \Delta$ में कोई $m$-स्थानी !!{predicate} है,
तो $\Assign{R}{M} = \Setabs{\tuple{t_1, \dots, t_m} \in
\Domain{M}^n}{R(t_1, \dots, t_n) \in \Theta}$।
\end{enumerate}
$\Struct{M}$ इस अर्थ में \emph{पद निदर्श} है कि उसका प्रांत पदों का समुच्चय
है, और !!{constant}s तथा !!{function}s की व्याख्या इस प्रकार की जाती है कि
$\Value{t}{M} = t$ सुनिश्चित हो।~$\Theta$ के परमाण्विक !!{formula}s का उपयोग
$\Assign{R}{M}$ को इस तरह परिभाषित करने में होता है कि
$\Sat{M}{R(t_1, \dots, t_n)}$ तभी और केवल तभी हो जब $R(t_1, \dots, t_n)
\in \Theta$।

\begin{lem}\ollabel{lem:truth}
प्रत्येक $!A \in \Theta \cup \Xi$ के लिए:
\begin{enumerate}
  \item यदि $!A \in \Theta$, तो $\Sat{M}{!A}$।
  \item यदि $!A \in \Xi$, तो $\Sat/{M}{!A}$।
\end{enumerate}
\end{lem}

\begin{proof}
  $\depth{!A}$ पर आगमन से।

पहले मान लें कि $!A$ परमाण्विक है, अर्थात् या तो $!A \ident \lfalse$ या $!A
\ident R(t_1, \dots, t_m)$।

चूँकि $\Pi_n \Sequent \Lambda_n$ विफलता शाखा है, कोई $\Pi_n$, $\lfalse$ को
समाहित नहीं कर सकता (अन्यथा $\Pi_n \Sequent \Lambda_n$ अभिगृहीत होता)।
\Struct{M} की परिभाषा से $\Sat{M}{R(t_1, \dots, t_m)}$ तभी और केवल तभी है
जब $R(t_1, \dots, t_m) \in \Theta$। दोनों को मिलाकर: यदि $!A \in \Theta$,
तो $\Sat{M}{!A}$।

यदि $!A$ परमाण्विक है और $!A \in \Pi_n$, तो $!A \in \Pi_{n+1}$; और यदि
$!A \in \Lambda_n$, तो $!A \in \Lambda_{n+1}$। (यह प्रमाण-खोज कलन-विधि
द्वारा उत्तराधिकारी अनुक्रम उत्पन्न करने के ढंग से मिलता है।) अतः यदि $!A
\in \Theta \cap \Xi$, तो किसी~$n$ के लिए $!A \in \Pi_n \cap \Lambda_n$।
इसका अर्थ है कि $\Pi_n \Sequent \Lambda_n$ अभिगृहीत है, जबकि विफलता शाखा में
कोई अभिगृहीत नहीं हो सकता। इसलिए रचना से परमाण्विक~$!A$ के लिए $!A \notin
\Theta \cap \Xi$; फलतः यदि $!A \in \Xi$, तो $!A \notin \Theta$।
\Struct{M} की परिभाषा से इसका अर्थ है कि यदि $!A \in \Xi$, तो
$\Sat/{M}{!A}$।

आगमन चरण के लिए मान लें कि (1) और (2),~$!A$ से कम !!{depth} वाले सभी
!!{formula}s के लिए सत्य हैं। स्थितियों में भेद करके हम $!A$ के लिए (1) और
(2) सिद्ध करते हैं।

पहले ध्यान दें कि जब $!A^i$, $\Pi_n \Sequent \Lambda_n$ में आता है, तब वह
$\Pi_{n+1} \Sequent \Lambda_{n+1}$ में भी आता है, जब तक कि $i$, $\Pi_n
\Sequent \Lambda_n$ का सबसे छोटा सूचकांक न हो। प्रत्येक चरण में जोड़े गए
सर्वोच्च अनुक्रमों में सबसे छोटा सूचकांक बढ़ता है, इसलिए अंततः हम ऐसे $\Pi_n
\Sequent \Lambda_n$ पर पहुँचते हैं जिसमें $!A^i$ आता है और $i$ सबसे छोटा
सूचकांक है। चरण~$n+1$ में कलन-विधि~$!A^i$ का अपचयन करती है। दूसरे शब्दों में,
यदि $!A \in \Theta$, तो किसी $n$ के लिए $\Pi_n = \Pi_n', !A^i$ और $i$
सबसे छोटा सूचकांक है। और यदि $!A \in \Xi$, तो किसी $n$ के लिए $\Lambda_n =
\Lambda_n', !A^i$ और $i$ सबसे छोटा सूचकांक है।

\begin{enumerate}
  \item $!A \in \Theta$ और $!A \ident !B \land !C$। तब किसी $n$ के लिए
  $\Pi_n = \Pi_n', (!B \land !C)^i$। $\Pi_{n+1} \Sequent \Lambda_{n+1}$,
  \LeftR{\land} का संगत आधार-वाक्य है, अर्थात् $\Pi_{n+1} = \Pi_n', !B^k,
  !C^{k+1}$। फलतः $!B, !C \in \Theta$। आगमन परिकल्पना से $\Sat{M}{!B}$ और
  $\Sat{M}{!C}$। $\Sat{M}{}$ की परिभाषा से $\Sat{M}{!B \land !C}$।

  \item $!A \in \Xi$ और $!A \ident !B \land !C$। तब किसी $n$ के लिए
  $\Lambda_n = \Lambda_n', !B \land !C$। $\Pi_{n+1} \Sequent
  \Lambda_{n+1}$, निष्कर्ष $\Pi_n \Sequent \Lambda'_n, !B \land !C$ वाले
  \RightR{\land} के दो आधार-वाक्यों में से एक है; अर्थात् $\Lambda_{n+1} =
  \Lambda_n', !B^k$ या $\Lambda_{n+1} = \Lambda_n', !C^k$। फलतः या तो $!B
  \in \Xi$ या $!C \in \Xi$। आगमन परिकल्पना से या तो $\Sat/{M}{!B}$ या
  $\Sat/{M}{!C}$। $\Sat{M}{}$ की परिभाषा से $\Sat/{M}{!B \land !C}$।
  
  \item $!A \ident \lnot !B$, $!A \ident !B \lor !C$ और $!A \ident !B
  \lif !C$ वाली स्थितियाँ समान हैं और अभ्यास के लिए छोड़ी जाती हैं।

  \item $!A \in \Theta$ और $!A \ident \lexists[x][!B(x)]$। तब किसी $n$ के
  लिए $\Pi_n = \Pi_n', \lexists[x][!B(x)]^i$। $\Pi_{n+1} \Sequent
  \Lambda_{n+1}$, \LeftR{\lexists} का संगत आधार-वाक्य है, अर्थात् किसी~$c$
  के लिए $\Pi_{n+1} = \Pi_n', !B(c)^k$। फलतः $!B(c) \in \Theta$ और~$c \in
  C$। आगमन परिकल्पना से $\Sat{M}{!B(c)}$। $\Sat{M}{}$ की परिभाषा से
  $\Sat{M}{\lexists[x][!B(x)]}$।

  \item $!A \in \Xi$ और $!A \ident \lexists[x][!B(x)]$। तब किसी $n$ के
  लिए $\Lambda_n = \Lambda_n', \lexists[x][!B(x)]^i$। हर बार
  $\lexists[x][!B(x)]$ का अपचयन होने पर वह नए सूचकांक सहित आधार-वाक्य अनुक्रम
  के दाएँ पक्ष पर बना रहता है; अतः विफलता शाखा में दाईं ओर आने पर उसका अनंत
  बार अपचयन होता है। प्रत्येक पद $t \in \Domain{M}$ अंततः विफलता शाखा पर
  ``केवल~$C_n$ के अचरों वाला ऐसा पहला पद जिसे दाईं ओर $\lexists[x][!B(x)]$
  के किसी अपचयन में अभी प्रयुक्त नहीं किया गया'' बनता है। अतः प्रत्येक $t \in
  \Domain{M}$ के लिए किसी~$n$ पर $!B(t) \in \Lambda_n$, और फलतः $!B(t) \in
  \Xi$। आगमन परिकल्पना से प्रत्येक $t \in \Domain{M}$ के लिए
  $\Sat/{M}{!B(t)}$। $\Sat{M}{}$ की परिभाषा से
  $\Sat/{M}{\lexists[x][!B(x)]}$।

  \item $!A \ident \lforall[x][!B(x)]$ वाली स्थितियाँ समान हैं और अभ्यास के
  लिए छोड़ी जाती हैं।
\end{enumerate}
\end{proof}

\begin{cor}\ollabel{cor:complete}
\Log{G3c} की प्रमाण-खोज कलन-विधि पूर्ण है: यदि वह बीच में रुकती है या कभी
समाप्त नहीं होती, तो आरंभिक अनुक्रम~$\Gamma \Sequent \Delta$ वैध नहीं है और
इसलिए उसका कोई !!{proof} नहीं है।
\end{cor}

\begin{proof}
  यदि कलन-विधि बीच में रुकती है या कभी समाप्त नहीं होती, तो ऐसी विफलता शाखा
  है जिससे $\Struct{M}$ परिभाषित किया जा सकता है। \olref{lem:truth} से प्रत्येक
  $!A \in \Gamma$ के लिए $\Sat{M}{!A}$ और प्रत्येक $!A \in \Delta$ के लिए
  $\Sat/{M}{!A}$। (याद रखें कि $\Pi_0 = \Gamma$ और $\Lambda_0 = \Delta$,
  इसलिए $\Gamma \subseteq \Theta$ और $\Delta \subseteq \Xi$।) अतः $\Gamma
  \Sequent \Delta$ वैध नहीं है। \Log{G3c} की निर्दोषता से $\Gamma \Sequent
  \Delta$ का कोई !!{proof} नहीं है।
\end{proof}

\begin{cor}
\Log{G3c} पूर्ण है: यदि $\Gamma \Sequent \Delta$ वैध है, तो $\Log{G3c}
\Proves \Gamma \Sequent \Delta$।
\end{cor}

\begin{proof}
  हम प्रतिधनात्मक कथन सिद्ध करते हैं। मान लें $\Log{G3c} \Proves/ \Gamma
  \Sequent \Delta$। तब प्रमाण-खोज कलन-विधि को बीच में रुकना होगा या कभी समाप्त
  नहीं होना होगा, क्योंकि खोजने के लिए कोई !!{proof} नहीं है।
  \olref{cor:complete} से $\Gamma \Sequent \Delta$ वैध नहीं है।
\end{proof}


\end{document}
